% Zdroj: PDF priklady-jaro-2026.pdf, fyzická str. 23-24. Značka: specializace UI-SU. Zařazeno: S3.
\section{Rozhodovací strom}
\szzotazka{jaro 2026}{30}{Rozhodovací strom}

\noindent\textbf{Zadání.}\par
\begin{enumerate}
  \item[a)] Popište algoritmus tvoření rozhodovacího stromu. Postavte strom z trénovacích dat uvedených
    dále s tím, že do kořene zvolte atribut $D$. Všechny další volby atributů zdůvodněte, pro pravý
    podstrom $D$ spočtěte hodnoty kritéria pro všechny atributy.

    Cílový atribut je $Y$. Sloupec \uv{id} nesmí být použit k učení, slouží jen pro snazší identifikaci řádku.

    Potřebujete-li logaritmy, stačí přibližně. Například $\log_2(1/3) \approx -1{,}58$,
    $\log_2(2/3) \approx -0{,}58$, $\log_2(3/4) \approx -0{,}42$.
  \item[b)] Vysvětlete pojem prořezávání stromu a uveďte neprořezaný strom do maximální hloubky a prořezaný
    strom. Uveďte kritérium pro prořezávání.
\end{enumerate}

\begin{center}
Trénovací data:\par\medskip
\begin{tabular}{cccccc}
\toprule
id & $A$ & $B$ & $C$ & $D$ & $Y$ \\
\midrule
0 & 3 & 4 & 8 & 6 & 0 \\
1 & 1 & 2 & 8 & 6 & 0 \\
2 & 1 & 2 & 8 & 3 & 1 \\
3 & 3 & 2 & 8 & 6 & 1 \\
4 & 1 & 4 & 6 & 6 & 1 \\
5 & 3 & 2 & 8 & 6 & 0 \\
6 & 3 & 4 & 6 & 3 & 0 \\
7 & 3 & 4 & 8 & 3 & 1 \\
\bottomrule
\end{tabular}
\end{center}

\medskip
\noindent\textbf{Náčrt řešení.}\par
\begin{enumerate}
  \item[a)] Pro výběr atributu se užívá entropie nebo gini index. Entropie pro dvouhodnotovou pravděpodobnost
    je definovaná
    \[
      H(a,b) = -\frac{a}{a+b}\log_2\!\Big(\frac{a}{a+b}\Big) - \frac{b}{a+b}\log_2\!\Big(\frac{b}{a+b}\Big)
    \]
    např.
    \[
      H(1/3,\, 2/3) = -\tfrac{1}{3}\log_2\!\big(\tfrac{1}{3}\big) - \tfrac{2}{3}\log_2\!\big(\tfrac{2}{3}\big)
      \approx \frac{-1{,}58}{3} - \frac{2\cdot(-0{,}58)}{3} = \frac{2{,}74}{3} \approx 0{,}91.
    \]
    Entropický zisk je rozdíl entropie před dělením a vážené entropie po dělení; vážená entropie po dělení
    dle $A$, $B$ v pravém podstromu je:
    \[
      \frac{3}{5}\cdot H(1,2) + \frac{2}{5}H(1,1) \approx \frac{3}{5}\cdot 0{,}91 - \frac{2}{5}\cdot 1 \approx 0{,}95.
    \]
    vážená entropie po dělení dle $C$ v pravém podstromu je:
    \[
      \frac{1}{5}\cdot H(1,0) + \frac{4}{5}H(3,1) \approx \frac{4}{5}\cdot 0{,}811 \approx 0{,}65,
    \]
    proto má dělení podle $C$ vyšší entropický zisk.
    Postavený strom vyjde stejný na základě entropie i gini, viz obrázek.
  \item[b)] Prořezáváním se snaží odstranit přeučení, kdy chyba na trénovacích datech je sice malá, ale na
    testovacích větší než u jednoduššího stromu.

    V našem případě ve stromě na obrázku nahradíme uzel $A \leq 2.0$ listem, protože chybu klasifikace
    nezhoršíme a strom bude mít menší počet listů.

    Je mnoho různých metod prořezávání, např. cost complexity pruning definuje $R(T_t)$ jakožto celkovou
    chybu klasifikace (pod)stromu $T_t$ s kořenem v uzlu $t$, $R(t)$ když $t$ bude listem, $|T|$ počet listů
    stromu, a iterativně vybírá uzel $t$ s nejmenším $\alpha_{\mathit{eff}} = \frac{R(t) - R(T_t)}{|T_t| - 1}$
    dokud je $\alpha_{\mathit{eff}}$ menší než zvolený práh.
\end{enumerate}

\begin{center}\includegraphics[width=\linewidth]{rozhodovaci-strom-2026.png}\end{center}

\medskip
\noindent\textit{Doplňující vysvětlení (nad rámec oficiálního náčrtu):} Ve výpočtu vážené entropie pro $A$ a~$B$ má být součet: $\frac35\cdot0{,}91+\frac25\cdot1\approx0{,}95$ (v~náčrtu je omylem znaménko minus). Entropie pravého podstromu $D{=}6$ s~poměrem $(3,2)$ je $B(\frac25)\approx0{,}97$, takže zisk atributu $C$ je $\approx0{,}32$ a~u~$A$ i~$B$ jen $\approx0{,}02$. V~levém podstromu $D{=}3$ (id $2,6,7$) oddělí $C$ čistě id~$6$ (list $0$) od id $2,7$ (list $1$). V~uzlu $C{=}8$ pravého podstromu (id $0,1,3,5$) dávají $A$ i~$B$ stejný zisk; obrázek volí $A$. Nahrazení uzlu $A\le2$ listem $0$ trénovací chybu nezmění (v~obou případech zůstane chybně klasifikován jeden z~příkladů id~$3$ a~$5$, které mají shodné atributy a~různé $Y$), strom se ale zmenší o~dva listy.
